\chapter{Исследование проблемной области "Процесс диагностики заболеваний дынь"}
\label{chapter1}

\section{Краткая характеристика источников знаний и методов получения знаний}

Для моделирования понятийной структуры проблемной области используются три основные категории источников знаний:
\begin{itemize}
    \item люди и эксперты, обладающие глубокими знаниями и опытом в соответствующей области;
    \item разнообразные литературные источники, включая книги, справочники, инструкции и прочие печатные материалы;
    \item базы данных, электронные носители и иные цифровые ресурсы.
\end{itemize}
В рамках данной работы в качестве источника знаний использовалась литература (источник второго типа). В частности, основой послужили труды профессора, доктора технических наук Г.В. Рыбиной: книга «Теория и технология построения интегрированных экспертных систем» и совместное издание с С.С. Параджановым «Технология построения динамических интеллектуальных систем». Эти материалы применялись непосредственно при разработке и конструировании прототипа динамической интегрированной экспертной системы.  Для тематической части- книга Кярыма Маметгулова, Бахрии Юсубовой ”Болезни бахчевых культур и меры борьбы против них”

\section{Глоссарий предметной области}

\begin{description}
    \item[Сад] участок земли с растениями, выращиваемыми в открытом грунте или защищённом пространстве для получения урожая.
    \item[Датчик] устройство для преобразования параметров окружающей среды в измеримые сигналы, пригодные для анализа и обработки информации.
    \item[Воздух] газовая смесь, в составе которой преобладают азот и кислород, образующие земную атмосферу.
    \item[Благоприятные условия воздуха] такой состав атмосферы, при котором растения могут длительное время развиваться и функционировать оптимально.
    \item[Температурный режим воздуха] количественное измерение степени нагрева или охлаждения атмосферного воздуха.
    \item[Влажность атмосферы] степень присутствия водяных паров в воздухе, влияющая на жизнедеятельность растений.
    \item[Почва] верхний плодородный слой земной коры, состоящий из органических и минеральных компонентов, жидкостей и газов, обеспечивающих развитие растительных организмов.
    \item[Влажность грунта] процентное содержание воды в почве относительно её полностью высушенной массы или на единицу объёма.
    \item[Содержание азотных компонентов] уровень концентрации азота в почвенной среде, необходимый для растений.
    \item[Содержание калийных компонентов] количество калия в почве, влияющего на обменные процессы и устойчивость культур.
    \item[Природная среда] совокупность физических, химических и биологических факторов, в которых развиваются растения.
    \item[Специалист по растениеводству] лицо с профессиональной подготовкой и опытом в области выращивания и ухода за культурами.
    \item[Визуальная диагностика] метод определения заболеваний и дефицитов по внешним признакам и изменениям внешнего вида растений.
    \item[Микроэлементы и макроэлементы] химические вещества, необходимые растениям для полноценного роста и развития.
    \item[Симптомы заболевания] видимые изменения структуры и окраски растения, указывающие на наличие болезни или стресса.
    \item[Динамическое содержание питательных веществ] колебание уровня элементов питания в почве с течением времени и сезонов.
    \item[Дыня (Cucumis melo)] культура из семейства тыквенных, характеризующаяся мясистыми плодами и требующая специального ухода при выращивании.
    \item[Мучнистая роса] грибковое поражение надземных частей растения, проявляющееся образованием светлого налёта на листовых пластинках и стеблях.
    \item[Фузариозное увядание] болезнь, вызванная грибковым возбудителем, приводящая к потере тургора и отмиранию растительной ткани.
    \item[Антракноз] грибковая инфекция, характеризующаяся формированием некротических пятен тёмного цвета на листьях и плодах с последующей гнилью.
    \item[Вредоносные организмы] живые существа, способные причинять вред растительным культурам и служить переносчиками инфекций.
\end{description}

\section{Описание задач/подзадач, решаемых в данной предметной области}

В настоящее время проблемой своевременной диагностики и предотвращения заболеваний занимается специалист по защите растений. Для эффективного выявления болезней дынь и оценки состояния растений при изменяющихся условиях питания необходима полная и достоверная информация о визуальных симптомах заболеваний, соответствующих характеру поражения. Данная задача является сложной для специалиста и требует постоянного анализа состояния культур и внешних факторов, влияющих на их здоровье. Высокая вариабельность проявления симптомов в условиях нестабильного содержания питательных веществ в почве усложняет процесс идентификации болезней и выбора адекватных мер борьбы. Для успешного выполнения диагностических и профилактических работ необходимо поставить и решить следующие задачи.

\subsection{Задача мониторинга}

Система в реальном времени получает визуальные данные о состоянии листьев и плодов дынь, информацию о содержании питательных веществ в почве и климатических параметрах окружающей среды, фиксирует изменения окраски, появление пятен, деформации и других визуальных признаков, сигнализируя о возможном развитии заболевания или дефицита элементов питания.
Постановка задачи:
Дано: Данные о визуальных признаках растений дынь (цвет листьев, наличие пятен, состояние плодов), показатели содержания азота, фосфора, калия, кальция и магния в почве, метеорологические данные (температура, влажность, освещённость).
Требуется: Отслеживать текущее состояние дынь в условиях нестабильного содержания питательных веществ и оповещать о появлении подозрительных визуальных симптомов, указывающих на развитие заболевания.

\subsection{Задача диагностики}

Система анализирует визуальные признаки поражённых растений дынь (характер пятен, налёты, деформации, изменение окраски листьев и плодов), соотносит выявленные симптомы с известными болезнями и дефицитами питательных веществ, определяет вероятный диагноз и предлагает рекомендации по лечению и профилактике на основе анализа условий выращивания.
Постановка задачи:
Дано: Визуальные данные о симптомах поражения дынь (фотографии, описание признаков), информация о содержании макро- и микроэлементов в почве, условиях выращивания (температура, влажность, освещённость).
Требуется: Определить вид заболевания или тип дефицита питательных веществ и сформировать набор рекомендаций по устранению выявленной проблемы с учётом текущих условий окружающей среды.

\section{Описание внешнего мира - сада с дынями}

Внешний мир сада, в котором выращиваются дыни, включает в себя природные факторы (атмосферные осадки, температурные колебания, солнечное излучение, ветровые потоки) и антропогенные воздействия (применение удобрений, использование пестицидов, механическое воздействие). Их влияние обусловлено физическими, химическими и биологическими процессами, протекающими в почве и атмосфере.
Так, например, интенсивные осадки повышают влажность почвы и могут привести к вымыванию питательных веществ, что способствует их нестабильному содержанию в грунте. Недостаток осадков, напротив, снижает доступность микроэлементов и может вызвать дефицит питания у растений. Повышение температуры воздуха ускоряет развитие грибковых заболеваний, таких как мучнистая роса, в то время как перепады влажности создают благоприятные условия для распространения патогенов.
Содержание азота, фосфора и калия в почве непрерывно изменяется под воздействием погодных условий, жизнедеятельности микроорганизмов и процессов выщелачивания. Эти колебания напрямую влияют на визуальное состояние растений, проявляясь в изменении окраски листьев, деформации побегов и снижении устойчивости к болезням.
Схематично взаимодействие сада и внешнего мира отражено на схеме, где дыни выступают системой, подвергающейся влиянию переменчивых условий окружающей среды и требующей постоянного мониторинга и диагностики состояния.

\section{Системный анализ на применимость технологии систем, основанных на знаниях}
Диагностика состояния сложных систем, таких как бахчевые культуры, основывается на анализе качественных данных (символов), а не на простых числовых расчетах. Эксперты считают эту задачу нетривиальной из-за множества входных параметров от датчиков и отсутствия прямого алгоритмического решения. Тем не менее, масштабы задачи позволяют эффективно использовать компьютерные системы. Таким образом, разработка интегрированной экспертной системы для этой предметной области является целесообразной.

Внедрение такой системы позволит значительно сократить расходы на обучение персонала и уменьшить количество специалистов, контролирующих состояние культур, а также объём специализированного оборудования. Следовательно, создание интегрированной экспертной системы экономически оправдано.

Решение этой задачи требует интеллектуального подхода и рассуждений, а не механических действий. Специалисты в данной области обладают знаниями о методах работы и согласны в подходах к решению, что подтверждает возможность создания интегрированной экспертной системы.

Исходя из вышеизложенного, применение методов и технологий динамических интегрированных экспертных систем в данной предметной области является обоснованным, уместным и оправданным.

\section{Постановка задачи}
Основной целью данной работы является создание прототипа динамической интегрированной экспертной системы (ИЭС). Эта система должна визуализировать текущие условия на поле и формировать диагностические заключения о состоянии культур, чтобы помочь агрономам оперативно принимать обоснованные решения.
